\clearpage
\fancyhead[L]{\small JUNO.}
\hypertarget{native-263-lower}{}
\begin{center}
JUNO.

\rule{0.6\linewidth}{0.4pt}

Monatliche Correspondenz. Vol. X. pp. 552–555. December 1804.

\rule{0.6\linewidth}{0.4pt}
\end{center}

\noindent—\quad—\quad—\par

Unfortunately, the sky here has been constantly overcast since 6 November, so we have been unable to continue observing Juno or to supply Dr. \textsc{Gauss} with new data to refine the elements he had already calculated. We did, however, send him two observations of our own from 20 and 21 Oct. They differed by $1\frac{1}{2}$ minutes from his second set of elements, which was enough to prompt this astronomer, as tireless as he was perceptive, to calculate a new third set of elements for Juno at once. He sent them to us with the following remarks. “Since my last letter I have calculated the following new third set of elements. I used the observations of Juno on 20 and 21 Oct. that you kindly sent me and that already differed by $1\frac{1}{2}$ minutes from the second set of elements.

\begin{center}
\begingroup\fontsize{8}{10}\selectfont\setlength{\tabcolsep}{3pt}\renewcommand{\arraystretch}{1.1}
\begin{tabular}{lr}
Epoch 1804, 30 Sept., at $0^{\mathrm u}$ on the Seeberg meridian . . & $22^\circ\ 34'\ 48''$ \\
Daily mean motion . . . . . . . . . . . & $842''75$ \\
Aphelion . . . . . . . . . . . . . . . . & $233^\circ\ 56'\ 6''$ \\
Logarithm of the semimajor axis . . . . . . . & $0.426699$ \\
Eccentricity . . . . . . . . . . . . . . . . & $0.263182$ \\
Ascending node . . . . . . . . . . . . . & $171^\circ\ 0'\ 0''$ \\
Inclination of the orbit . . . . . . . . . . . . . . & $12\ 52\ 48$ \\
\end{tabular}\endgroup
\end{center}

“I have compared these elements with all your observations and found the following agreement.
\SourceLeaf{264}
\begin{center}
\begingroup\fontsize{8}{10}\selectfont\setlength{\tabcolsep}{3pt}\renewcommand{\arraystretch}{1.1}
\begin{tabular}{r|r|r|r|r|r}
1804 & \shortstack{Mean time\\at\\Seeberg} & \shortstack{Calculated right\\ascension\\of $\JunoSymbol$} & \shortstack{Calculated southern\\declination\\of $\JunoSymbol$} & \multicolumn{2}{c}{\shortstack{Difference\\R—B\\in R.A.\qquad in decl.}}\\ \hline
13 Sept. & $12^{\mathrm u}\ 31^{\mathrm m}\ 59^{\mathrm s}$ & $0^\circ\ 44'\ 55''8$ & $1^\circ\ 52'\ 34''3$ & $-0''6$ & $-2''7$ \\
14 & $12\ 27\ 27$ & $\phantom{000}\ 35\ 43.2$ & $2\ 5\ 35.3$ & $+2.2$ & $-0.2$ \\
15 & $12\ 22\ 53$ & $\phantom{000}\ 26\ 20.8$ & $\phantom{0}\ 18\ 42.2$ & $+0.8$ & $+3.7$ \\
17 & $12\ 13\ 45$ & $\phantom{000}\ 7\ 10.6$ & $\phantom{0}\ 45\ 9.2$ & $+2.5$ & $-3.9$ \\
18 & $12\ 9\ 10$ & $359\ 57\ 24.5$ & $\phantom{0}\ 58\ 28.2$ & $+4.0$ & $-2.7$ \\
20 & $.\ .\ .\ .$ & $\phantom{000}\ 37\ 35.7$ & $3\ 25\ 32.7$ & $.\ .\ .\ .$ & $+9.2$ \\
23 & $11\ 46\ 11$ & $\phantom{000}\ 7\ 25.9$ & $4\ 5\ 22.9$ & $+11.0$ & $-11.6$ \\
27 & $11\ 27\ 46$ & $358\ 27\ 7.7$ & $\phantom{0}\ 58\ 29.8$ & $+7.0$ & $-27.4$ \\
28 & $11\ 23\ 11$ & $\phantom{000}\ 17\ 9.6$ & $5\ 11\ 35.1$ & $+3.2$ & $-8.8$ \\
30 & $11\ 14\ 1$ & $357\ 57\ 28.6$ & $\phantom{0}\ 37\ 28.2$ & $-1.9$ & $-15.5$ \\
2 Oct. & $11\ 4\ 52$ & $\phantom{000}\ 38\ 15.5$ & $6\ 2\ 50.8$ & $+0.5$ & $-8.7$ \\
4 & $10\ 55\ 46$ & $\phantom{000}\ 19\ 40.9$ & $\phantom{0}\ 27\ 36.7$ & $+1.4$ & $-2.0$ \\
5 & $10\ 51\ 15$ & $\phantom{000}\ 10\ 41.1$ & $\phantom{0}\ 39\ 43.8$ & $-0.4$ & $+3.9$ \\
6 & $10\ 46\ 44$ & $\phantom{000}\ 1\ 54.0$ & $\phantom{0}\ 51\ 39.0$ & $-2.3$ & $+6.8$ \\
10 & $10\ 28\ 50$ & $356\ 29\ 14.9$ & $7\ 37\ 10.5$ & $+0.8$ & $+19.4$ \\
12 & $10\ 20\ 0$ & $\phantom{000}\ 14\ 37.8$ & $7\ 58\ 30.8$ & $-5.7$ & $+4.3$ \\
20 & $9\ 45\ 34$ & $355\ 29\ 55.3$ & $9\ 12\ 58.8$ & $+3.4$ & $-4.6$ \\
21 & $9\ 41\ 23$ & $\phantom{000}\ 26\ 2.7$ & $\phantom{0}\ 20\ 45.8$ & $+2.6$ & $+2.6$ \\
\end{tabular}\endgroup
\end{center}

“Dr. \textsc{Maskelyne} kindly sent me the following three additional observations.

\begin{center}
\begingroup\fontsize{9}{11}\selectfont\setlength{\tabcolsep}{3pt}\renewcommand{\arraystretch}{1.1}
\begin{tabular}{r|r|r|r}
1804 & \shortstack{Mean time\\at\\Greenwich} & \shortstack{Apparent right\\ascension\\of $\JunoSymbol$} & \shortstack{Apparent southern\\declination\\of $\JunoSymbol$}\\ \hline
5 Oct. & $10^{\mathrm u}\ 51^{\mathrm m}\ 6^{\mathrm s}$ & $357^\circ\ 10'\ 24''6$ & $6^\circ\ 40'\ 1''1$ \\
9 & $10\ 33\ 9$ & $356\ 36\ 46.9$ & $7\ 26\ 21.3$ \\
17 & $9\ 58\ 10$ & $355\ 43\ 45.6$ & $8\ 47\ 18.5$ \\
\end{tabular}\endgroup
\end{center}

“The third set of elements agrees with these and the earlier observations as shown below.

\begin{center}
\begingroup\fontsize{8}{10}\selectfont\setlength{\tabcolsep}{3pt}\renewcommand{\arraystretch}{1.1}
\begin{tabular}{r|r|r|r|r|r}
1804 & \shortstack{Mean time\\at\\Greenwich} & \shortstack{Calculated R.A.\\of $\JunoSymbol$} & \shortstack{Calculated southern\\declination\\of $\JunoSymbol$} & \multicolumn{2}{c}{\shortstack{Difference\\R—B\\in R.A.\qquad in decl.}}\\ \hline
25 Sept. & $11^{\mathrm u}\ 36^{\mathrm m}\ 50^{\mathrm s}$ & $358^\circ\ 46'\ 55''7$ & $4^\circ\ 32'\ 26''8$ & $+7''7$ & $-1''4$ \\
29 & $11\ 18\ 27$ & $358\ 6\ 58.7$ & $5\ 24\ 58.3$ & $+6.4$ & $+1.6$ \\
5 Oct. & $10\ 51\ 6$ & $357\ 10\ 24.7$ & $6\ 40\ 6.0$ & $+0.1$ & $+4.9$ \\
9 & $10\ 33\ 9$ & $356\ 36\ 46.4$ & $7\ 26\ 28.0$ & $-0.5$ & $+6.7$ \\
17 & $9\ 58\ 10$ & $355\ 43\ 46.0$ & $8\ 47\ 27.1$ & $+0.4$ & $+8.6$ \\
\end{tabular}\endgroup
\end{center}

“According to the third set of elements, the future course of \emph{Juno} is as follows. [From 2 Nov. 1804 to 4 Jan. 1805 at three-day intervals, right ascension, declination and brightness.]

“The unit of brightness is the brightness the planet would have at a distance of 1 from the Sun and from Earth. On that same scale its brightness was

\begin{center}
\begingroup\fontsize{9}{11}\selectfont\setlength{\tabcolsep}{3pt}\renewcommand{\arraystretch}{1.1}
\begin{tabular}{lrl}
on 5 Sept. & $0.1378$ &  \\
— 12 — & $0.1484$ &  \\
— 3 Oct. & $0.1640$ & highest this year \\
— 21 — & $0.1543$ &  \\
\end{tabular}\endgroup
\end{center}

“It will be interesting to see how long \emph{Juno} will remain visible this time. At the next opposition, in early March 1806, in the constellation Leo, it reaches only a quarter of its greatest brightness this year. For 31 December 1805, I find its position to be
\SourceLeaf{265}
\begin{center}
R.A. $176^\circ\ 45'$\quad Southern decl. $2^\circ\ 44'$\quad Brightness $0.0284$.
\end{center}

\Needspace{17\baselineskip}
\indent Here are all the observations we were able to obtain.

\begin{center}
\begingroup\fontsize{9}{11}\selectfont\setlength{\tabcolsep}{3pt}\renewcommand{\arraystretch}{1.1}
\begin{tabular}{r|r|r|r}
1804 & \shortstack{Mean time\\at\\Seeberg} & \shortstack{Apparent right\\ascension\\of $\JunoSymbol$} & \shortstack{Apparent southern\\declination\\of $\JunoSymbol$}\\ \hline
23 Oct. & $9^{\mathrm u}\ 33^{\mathrm m}\ 4^{\mathrm s}338$ & $355^\circ\ 19'\ 23''69$ & $9^\circ\ 35'\ 53''6$ \\
24 & $9\ 28\ 57.709$ & $355\ 16\ 42.51$ & $.\ .\ .\ .$ \\
30 & $9\ 4\ 54.644$ & $355\ 9\ 47.25$ & $10\ 18\ 43.8$ \\
5 Nov. & $8\ 41\ 56.099$ & $355\ 19\ 2.57$ & $.\ .\ .\ .$ \\
18 & $7\ 55\ 51.375$ & $356\ 34\ 46.58$ & $10\ 58\ 16.2$ \\
\end{tabular}\endgroup
\end{center}

\hfill[\textsc{von Zach.}]
